% !TEX root = ../main.tex
\chapter{Execució a bord i discussió}\label{ch:bord}
\fancyhead[RE]{\slshape \textsf{Capítol \thechapter.\ Execució a bord}}
\fancyhead[LO]{\slshape \textsf{\nouppercase \rightmark}}

\section{Per què cal avaluar el processament a bord}\label{sec:onboard_motivation}

La utilitat del control espacial de qualitat depèn de poder prendre la decisió
abans de transmetre les dades. En una estació terrestre seria possible aplicar
un segmentador més gran, revisar manualment la imatge o repetir la compressió.
En canvi, a bord hi ha restriccions de còmput, memòria, energia i temps de
contacte. Per aquest motiu, no és suficient demostrar que una màscara produeix
menys bitrate: també cal mesurar el cost de generar-la i d'executar el còdec.

L'arquitectura escollida separa dos costos:

\begin{enumerate}[leftmargin=*]
  \item \textbf{decisió espacial}: lectura del RAW, càlcul de l'índex o
  descriptor del preset i escriptura del Q-map;
  \item \textbf{transformació neuronal}: anàlisi espectral i espacial,
  modulació, quantització i generació dels latents.
\end{enumerate}

Aquesta separació és rellevant des de la perspectiva de sistema. Si el primer
cost fos comparable al segon, l'estalvi de bits podria no compensar el càlcul
addicional. Si és molt menor, es poden oferir diversos presets sense alterar
de manera apreciable la latència del còdec.

\section{Plataforma i abast de la demostració}\label{sec:rpi_platform}

La plataforma és una Raspberry Pi 3 Model B+ amb quatre nuclis ARM
Cortex-A53 a 1,4~GHz i aproximadament 1~GiB de memòria. El sistema utilitza
Linux de 64 bits i compila el codi amb GCC, \code{-O3}, OpenMP i opcions
específiques per a Cortex-A53. Aquesta família de processadors és pertinent per
estudiar codi de baix recurs i comparteix característiques amb ordinadors de
bord basats en ARM, però la Raspberry no és maquinari qualificat per a vol.
Aspectes com radiació, redundància, temps real, consum elèctric, interfícies del
sensor i certificació queden fora de la demostració. La comparació correcta és,
per tant, de \emph{representativitat computacional}, no d'equivalència de vol
\cite{raspberry2018brief,schwaller2017a53,serra2024c3satp}.

El bundle d'assaig és autocontingut: inclou codi C, pesos, calibració, sis
retalls RAW, llindars i scripts. No instal·la TensorFlow, NumPy ni paquets del
sistema. Cada execució desa \code{/usr/bin/time -v}, RSS màxim, temperatura,
estat de throttling, arguments, hashes i mides. Aquesta decisió limita les
diferències entre l'entorn local i la placa i permet reconstruir posteriorment
les mètriques al PC.

\begin{table}[H]
\centering
\small
\begin{tabular}{ll}
\toprule
Element & Configuració mesurada \\
\midrule
CPU & 4 $\times$ ARM Cortex-A53, fins a 1,4~GHz \\
Memòria & 909~MiB visibles pel sistema \\
Entrada & 8 bandes, $512\times512$, \code{uint16}, 4.194.304 bytes \\
Compilació & GCC \code{-O3 -fopenmp -mcpu=cortex-a53} \\
Paral·lelisme final & 4 fils OpenMP \\
Mostra & 6 retalls $\times$ 11 polítiques = 66 execucions \\
Instrumentació & \code{/usr/bin/time -v}, \code{vcgencmd}, hashes i logs \\
\bottomrule
\end{tabular}
\caption{Condicions principals del benchmark Raspberry optimitzat. La mostra
correspon a sis retalls i onze polítiques amb quatre fils.}
\label{tab:rpi_setup}
\end{table}

\section{Escalat inicial amb OpenMP}\label{sec:rpi_scaling}

Abans de modificar els kernels es va executar la mateixa càrrega amb un, dos i
quatre fils. Els valors de la Taula~\ref{tab:rpi_scaling} són mitjanes per una
imatge i una política; no són el temps de tot el dataset. Cada entrada conté
2.097.152 mostres espectrals.

\begin{table}[H]
\centering
\small
\begin{tabular}{rrrrr}
\toprule
Fils & Compressió (s) & Descompressió (s) & Total (s) & Speedup total \\
\midrule
1 & 561,73 & 413,22 & 974,96 & 1,00$\times$ \\
2 & 368,20 & 270,85 & 639,05 & 1,53$\times$ \\
4 & 316,43 & 247,69 & 564,12 & 1,73$\times$ \\
\bottomrule
\end{tabular}
\caption{Escalat inicial mitjà per retall i política. El speedup es calcula
respecte d'un fil.}
\label{tab:rpi_scaling}
\end{table}

Quatre fils redueixen el temps total un 42,1\% respecte d'un fil. L'escalat no
és lineal: un speedup ideal seria $4\times$, però s'obté $1,73\times$. Hi
contribueixen les regions serials, el cost de crear equips OpenMP, la pressió
sobre memòria i l'amplada de banda compartida. Tot i això, quatre fils és la
configuració operativa perquè és la més ràpida de les tres i no produeix
errors ni falta de memòria.

\subsection{Lectura amb la llei d'Amdahl}

La llei d'Amdahl no és una hipòtesi del projecte, sinó una eina clàssica per
interpretar per què un programa no accelera linealment quan s'afegeixen nuclis
\cite{amdahl1967}. En aquest cas serveix per donar context al resultat OpenMP:
si una part del temps no escala amb els fils, el speedup total queda limitat
encara que algunes funcions internes sí siguin paral·leles. La forma utilitzada
expressa el speedup amb $p$ fils com:

\begin{equation}
S(p)=\frac{1}{f+(1-f)/p},
\end{equation}

on $f$ és la fracció serial o, més precisament, la fracció aparent que no escala
en aquesta mesura. Substituint $S(4)=974,96/564,12=1,728$, s'obté:

\begin{equation}
f\simeq
\frac{1/S(4)-1/4}{1-1/4}=0,44.
\end{equation}

Aquest 44\% no és una mesura directa d'una funció serial del codi. Agrupa I/O,
regions no paral·lelitzades, ample de banda de memòria, sincronització, cost
d'OpenMP i possibles variacions de freqüència. La lectura útil és que afegir més
fils al mateix codi no pot donar un speedup proper al nombre de nuclis. Cal
reduir el treball de les capes o canviar la jerarquia de memòria.

La millora incremental ataca part d'aquesta fracció en GDN/IGDN,
depth-to-space i transformada espectral. Les convolucions continuen sent el
volum principal d'operacions, però ja disposaven de paral·lelisme per canals.
Una optimització posterior hauria de perfilar cache misses, vectorització i
reutilització de kernel abans d'afegir directives OpenMP indiscriminadament.

\section{Optimització incremental del C}\label{sec:c_optimization}

L'optimització no es va fer canviant el model, els pesos o l'ordre aritmètic
intern. Es van introduir tres canvis separats i es va executar una auditoria de
paritat després de cadascun:

\begin{description}[leftmargin=3.7cm,style=nextline]
  \item[GDN i IGDN] Paral·lelisme sobre píxels, mantenint privat el denominador
  de cada fil i l'ordre d'acumulació sobre canals.
  \item[Depth-to-space] Paral·lelisme sobre canals de sortida. Aquesta operació
  només reordena memòria i ha de ser exacta bit a bit.
  \item[Transformada espectral] Paral·lelisme sobre píxels, mantenint l'ordre
  intern de les bandes per no alterar l'aritmètica en coma flotant.
\end{description}

Abans dels canvis es va crear una golden baseline amb Q-map, bitstream,
reconstrucció, hashes i temps. L'auditor compara Q-map byte a byte i
reconstrucció RAW byte a byte. Si aparegués una diferència, també calcularia
error màxim, MAE, MSE i PSNR. En les 66 parelles Raspberry de la comparació
final no es va registrar cap fallada de paritat.

\subsection{Temps abans i després}

\begin{table}[H]
\centering
\small
\begin{tabular}{lrrrr}
\toprule
Etapa & Baseline (s) & Optimitzat (s) & Reducció (s) & Speedup \\
\midrule
Compressió & 316,43 & 261,32 & 55,11 & 1,211$\times$ \\
Descompressió & 247,69 & 206,80 & 40,89 & 1,198$\times$ \\
Total & 564,12 & 468,13 & 96,00 & 1,205$\times$ \\
\bottomrule
\end{tabular}
\caption{Temps mitjans per retall i política amb quatre fils. Les 66 parelles
comparen les mateixes entrades i Q-maps.}
\label{tab:rpi_optimized_time}
\end{table}

La millora total mitjana és $1,205\times$, equivalent a una reducció del
17,0\% del temps. La Figura~\ref{fig:rpi_speedup} mostra que el resultat és
estable entre polítiques: els speedups totals se situen aproximadament entre
$1,20\times$ i $1,21\times$. Això és coherent amb el fet que les polítiques
només canvien el Q-map; les capes convolucionals dominants són les mateixes.

\begin{figure}[H]
\centering
\includegraphics[width=0.88\textwidth]{raspberry_speedup_by_mode.png}
\caption{Speedup total de la versió C optimitzada per política, sempre amb
quatre fils. La línia $1\times$ representa absència de millora.}
\label{fig:rpi_speedup}
\end{figure}

El RSS màxim mitjà del compressor passa de 90.136~KiB a 89.505~KiB. Per tant,
el paral·lelisme introduït no exigeix una còpia completa del tensor per fil i
no augmenta la memòria de manera rellevant. Aquesta propietat és tan important
com el speedup en una placa amb menys d'1~GiB disponible.

\subsection{Paritat funcional}

La comparació revisa Q mínima, màxima, mitjana i nombre de valors únics, PSNR
global, PSNR ROI, PSNR de fons i proxy gzip. Es van emparellar 66 execucions i
es van obtenir zero fallades. El resultat accepta H5: les optimitzacions
redueixen temps sense canviar la decisió de qualitat ni la reconstrucció
validada.

La paritat amb la versió C anterior és el criteri d'acceptació estricte. La
comparació amb la referència Python de SORTENY continua essent una comprovació
externa del model, però no s'utilitza per atribuir a l'optimització diferències
que ja existien entre TensorFlow i la implementació C.

\section{Cost de generar el Q-map}\label{sec:qmap_cost}

La prova específica executa únicament els generadors de Q-map. Utilitza sis
entrades, onze casos i tres repeticions, és a dir, 198 mesures. No executa ni
el compressor ni el descompressor. Això aïlla el cost auxiliar de seleccionar
ROI, consultar la informació de qualitat i escriure el Q-map.

Aquest cost correspon al moment auxiliar de decisió: avaluar el preset sobre la
imatge, consultar la informació de qualitat disponible i escriure 1.024 valors
Q. No inclou la passada de còdec necessària per obtenir error base ni aplicar el
Q-map final. Aquesta separació és la que permet comprovar si la intel·ligència
espacial afegida és cara o si el coll d'ampolla continua sent el còdec.

\begin{table}[H]
\centering
\small
\begin{tabular}{lrrrr}
\toprule
Tipus & Mesures & Temps mitjà (s) & p95 (s) & RSS mitjà (KiB) \\
\midrule
\code{q204} uniforme & 18 & 0,066 & 0,070 & 1.580 \\
Adaptive-s8 & 18 & 0,067 & 0,070 & 1.614 \\
Presets semàntics & 162 & 0,113 & 0,120 & 5.459 \\
Tots els generadors & 198 & 0,104 & -- & -- \\
\bottomrule
\end{tabular}
\caption{Cost a Raspberry dels generadors de Q-map sobre sis retalls, onze
casos i tres repeticions. \emph{Mesures} és el nombre d'execucions
temporitzades. El percentil 95 (p95) indica que el 95\% de mesures no supera
el temps mostrat.}
\label{tab:qmap_cost}
\end{table}

El preset semàntic mitjà representa el 0,043\% del temps de compressió i el
0,024\% del temps de compressió més descompressió. En altres paraules, el
compressor tarda aproximadament 2.320 vegades més que la decisió semàntica
mitjana. Fins i tot \code{cloud\_avoid}, el cas més lent amb 0,173~s de
mitjana, queda per sota del 0,1\% del compressor.

Aquesta comparació separa dues parts del problema. Detectar la ROI o consultar
una taula auxiliar és molt barat. El cost fort de la ruta fixed-quality amb error
mesurat és obtenir aquest error amb el còdec real: una passada de compressor
costa 261,32~s per retall i compressor més descompressor costa 468,13~s. Per
tant, les optimitzacions C i el benchmark Raspberry no són només una validació
posterior; quantifiquen el cost de calcular $\lambda$/Q a bord quan la decisió
depèn de l'error real de la imatge. La conclusió és que el detector no és el
coll d'ampolla: el factor limitant és el còdec que proporciona l'error base i
aplica el Q-map.

\begin{figure}[H]
\centering
\includegraphics[width=0.88\textwidth]{raspberry_qmap_cost.png}
\caption{Temps mitjà de generació del Q-map comparat amb compressor i pipeline
complet. S'utilitza escala logarítmica perquè hi ha més de tres ordres de
magnitud de diferència.}
\label{fig:qmap_cost}
\end{figure}

Aquest resultat accepta H6 dins la mida d'entrada mesurada: el cost de decidir
la ROI automàtica és negligible davant del cost de les transformades
neuronals. No implica que qualsevol segmentador sigui gratuït. La conclusió és
específica dels presets algebraics implementats, que recorren les bandes i
calculen estadístics simples per bloc.

\section{Latència i throughput}\label{sec:throughput}

Un retall conté:

\begin{equation}
512\cdot512=262.144\;\text{posicions espacials},
\qquad
8\cdot512\cdot512=2.097.152\;\text{mostres espectrals}.
\end{equation}

Amb el compressor optimitzat, el throughput observat és:

\begin{align}
T_{espacial,enc} &= \frac{262.144}{261,32}
                   = 1.003\;\text{píxels espacials/s},\\
T_{mostres,enc} &= \frac{2.097.152}{261,32}
                   = 8.025\;\text{mostres espectrals/s}.
\end{align}

Quan s'inclou la descompressió, el temps total és 468,13~s i els valors baixen
a:

\begin{align}
T_{espacial,enc+dec} &= \frac{262.144}{468,13}
                   = 560\;\text{píxels espacials/s},\\
T_{mostres,enc+dec} &= \frac{2.097.152}{468,13}
                   = 4.480\;\text{mostres espectrals/s}.
\end{align}

La Taula~\ref{tab:throughput} evita confondre els dos escenaris.

\begin{table}[H]
\centering
\small
\begin{tabularx}{\textwidth}{lrrr>{\raggedright\arraybackslash}X}
\toprule
Ruta & Temps/retall & Píxels/s & Mostres/s & Ús principal \\
\midrule
Només compressor & 261,32~s & 1.003 & 8.025 & Codificació abans del downlink \\
Compressor+descompressor & 468,13~s & 560 & 4.480 & Càlcul d'error base o reconstrucció local \\
\bottomrule
\end{tabularx}
\caption{Throughput mesurat per un retall $8\times512\times512$. ``Píxel''
compta una posició espacial; ``mostra'' compta cada banda per separat.}
\label{tab:throughput}
\end{table}

La latència de compressió és 4~min 21~s per retall. Aquest valor és compatible
amb una demostració funcional i amb processament diferit de poques adquisicions,
però no permet afirmar per si sol processament en temps real. Una afirmació de
temps real exigiria comparar aquest throughput amb el cabal documentat del
sensor, la memòria de buffer i el cicle de servei de la missió concreta.

Com a ordre de magnitud extern, ESA descriu Sentinel-2 com una missió
multiespectral d'alt cabal que pot adquirir, emmagatzemar i descarregar fins a
1,6~TB per òrbita \cite{esa_sentinel2_operations}. Aquesta
referència no és un requisit per a aquest treball, perquè el sensor, el
maquinari i la missió són diferents, però ajuda a interpretar el resultat: una
Raspberry processant 55~MiB/h de dades crues no és una solució de temps real per
a sensors d'alt flux. En canvi, sí pot ser coherent amb processament diferit,
baixa cadència o selecció prèvia de pocs retalls. El resultat també és
consistent amb la motivació de plataformes a bord més especialitzades, com les
demostracions C3SatP/MENUT en òrbita \cite{serra2024c3satp}.

En la ruta fixed-quality amb error mesurat, el descompressor és rellevant perquè
permet obtenir la distorsió real de la passada base. En una ruta mínima de
transmissió amb el Q-map decidit sense error base, en canvi, el satèl·lit podria
executar només generació de Q-map més compressor i deixar la reconstrucció a
terra. La Taula~\ref{tab:operational_routes_cost} resumeix aquesta diferència.

\begin{table}[H]
\centering
\footnotesize
\begin{tabularx}{\textwidth}{>{\raggedright\arraybackslash}p{0.32\textwidth}r>{\raggedright\arraybackslash}X}
\toprule
Ruta & Temps & Lectura \\
\midrule
Q-map sense error base + compressió final & 261,4~s &
Cost mínim quan la qualitat local es decideix sense mesurar error base. \\
Error mesurat + Q-map + compressió final & 729,6~s &
Mesura error base, decideix qualitat i codifica el resultat final. \\
Taula precalibrada + Q-map + compressió final & 261,4~s &
Reutilitza informació calibrada a terra i evita la passada base a bord. \\
\bottomrule
\end{tabularx}
\caption{Cost per retall de tres rutes operatives sobre Raspberry. Els temps
combinen el Q-map semàntic mitjà (0,113~s) amb les mitjanes optimitzades de la
Taula~\ref{tab:rpi_optimized_time}.}
\label{tab:operational_routes_cost}
\end{table}

Els càlculs de la taula són:

\begin{equation}
0,113+261,32=261,43~\text{s},
\qquad
468,13+0,113+261,32=729,56~\text{s}.
\end{equation}

La taula precalibrada és important perquè desplaça a terra una part cara de la
decisió. No substitueix la validació de qualitat ni és universal: només és
fiable si la calibració representa bé el sensor, el tipus d'escena i les zones
que es volen protegir. Però, quan aquesta condició es compleix, pot ser una ruta
molt atractiva per a missió: manté el cost a bord pràcticament igual que una
compressió directa i evita aproximadament 468~s per retall respecte del càlcul
amb error mesurat. Per això es tracta com una opció operativa rellevant, no com
un simple detall auxiliar.

\subsection{Capacitat de servei per retalls}

Una altra manera d'interpretar la latència és com una cua de treball. Si la
Raspberry dedica un nucli de sistema i quatre fils OpenMP al compressor sense
interrupcions, la taxa màxima observada és:

\begin{equation}
\mu=\frac{3600}{261,32}=13,77\;\text{retalls/hora}.
\end{equation}

Això correspon aproximadament a 330 retalls per dia en un escenari ideal de
servei continu. Com cada RAW ocupa 4~MiB, el processador absorbiria uns
55~MiB/h de dades crues. Aquests valors són una capacitat de laboratori, no un
pressupost de missió: no inclouen adquisició, escriptura a emmagatzematge,
telemetria, altres tasques de l'OBC, períodes de baixa potència ni reinicis.

Si el sensor produeix retalls a una taxa mitjana $\lambda_a$, la cua només és
estable si:

\begin{equation}
\lambda_a<\mu.
\end{equation}

Una adquisició en ràfega podria superar temporalment $\mu$ si existeix buffer,
però el backlog hauria de buidar-se abans de la següent campanya. Aquesta
formulació és més útil que afirmar simplement que quatre minuts són molts:
permet inserir en el futur el nombre real de retalls per òrbita i comprovar si
la latència és compatible.

\subsection{Què falta per a un throughput de missió}

Per convertir la capacitat de laboratori en una conclusió de sistema calen:

\begin{itemize}[leftmargin=*]
  \item píxels o mostres generades per segon pel sensor;
  \item durada i freqüència de les adquisicions;
  \item memòria disponible per absorbir ràfegues;
  \item temps de CPU reservat per al compressor;
  \item potència mitjana i energia per retall;
  \item cabal i finestres de downlink.
\end{itemize}

Sense aquests paràmetres no es pot decidir si el 8,84\% de bitrate allibera més
recurs de sistema del que consumeix el càlcul complet. El projecte sí resol una
part concreta: el detector semàntic només representa el 0,043\% del temps de
compressió. La incertesa operativa resideix en el cost del còdec, en l'energia
per retall i en el perfil de missió.

\section{Temperatura, alimentació i repetibilitat}\label{sec:thermal}

Els registres de Raspberry indiquen valors de \code{get\_throttled} diferents
de zero, incloent \code{0xd0005} i \code{0xd0008}, i màxims de temperatura pròxims als
70~\textdegree C en algunes execucions. Això indica que les mesures es van fer
amb incidències històriques o actives d'alimentació i limitació de freqüència.
Els temps són reals de la placa utilitzada, però no representen una cota ideal
del Cortex-A53 amb alimentació estabilitzada i refrigeració controlada.

La comparació optimitzat--baseline continua sent útil perquè les dues campanyes
utilitzen la mateixa placa, el mateix nombre de fils i la mateixa càrrega.
Tanmateix, una campanya de caracterització de hardware formal hauria de fixar
governor, font d'alimentació, temperatura ambient i període de refredament, i
hauria de mesurar energia per imatge, no només temps.

\section{Cost computacional enfront d'estalvi de bits}\label{sec:cost_benefit}

La decisió final combina els Capítols~\ref{ch:resultats} i~\ref{ch:bord}. En la
mostra CSMR principal, \code{focus\_bgq128} estalvia un 8,84\% de bitrate
respecte de \code{q204}. A Raspberry, calcular el preset semàntic afegeix de
mitjana 0,113~s, un 0,043\% del temps del compressor. Per tant, dins del
pipeline actual, \textbf{sí que compensa executar el preset}. La pregunta més
difícil és si compensa la ruta completa de càlcul de qualitat amb error mesurat:
una compressió individual costa 261,32~s i compressor més descompressor costa
468,13~s per retall. Aquest cost pot ser acceptable en escenaris de baixa
cadència o processament diferit, però exigeix més optimització o hardware més
potent si es vol operar amb un flux d'adquisició alt.

Aquesta conclusió no significa que el sistema complet sigui prou ràpid. El
coll d'ampolla és el còdec neuronal C, no la selecció ROI. El treball futur
hauria de prioritzar kernels, vectorització, acceleració o un processador més
potent abans de simplificar els presets. El guany de compressió és moderat i
no justifica qualsevol cost computacional; justifica el cost concret mesurat
dels detectors algebraics.

\section{Discussió de viabilitat}\label{sec:onboard_discussion}

Els resultats permeten separar quatre nivells de viabilitat. Aquesta separació
evita barrejar una demostració funcional amb una certificació de missió:

\begin{table}[H]
\centering
\footnotesize
\begin{tabularx}{\textwidth}{>{\raggedright\arraybackslash}p{0.22\textwidth}>{\raggedright\arraybackslash}X}
\toprule
Nivell & Lectura \\
\midrule
Funcional &
Assolit: el bundle C executa compressor, descompressor i artefactes amb menys
d'1~GiB. Límit: no és software qualificat per a vol. \\
Decisió espacial &
Assolit: els presets automàtics costen 0,113~s de mitjana i no requereixen
operador. Límit: conclusió vàlida per als presets algebraics implementats. \\
Qualitat/bitrate &
Assolit dins l'abast: CSMR confirma estalvi de bitrate i ROI conservada en la
mostra seleccionada. Límit: dataset i màscares limitades. \\
Missió &
No demostrat: faltaria fixar cabal de sensor, energia, deadlines, tolerància a
fallades i hardware qualificat. \\
\bottomrule
\end{tabularx}
\caption{Lectura de viabilitat separada per nivells. Cada fila indica què queda
demostrat i quin límit impedeix convertir-ho directament en certificació de
vol.}
\label{tab:onboard_viability_levels}
\end{table}

La formulació evita dos extrems: presentar Raspberry com una certificació de
vol o descartar el resultat perquè no és el computador final. La placa serveix
per descobrir que la memòria és assumible, que quatre fils milloren el temps,
que els presets són pràcticament gratuïts i que el còdec encara necessita una
millora de throughput per a escenaris exigents.

En conjunt, H1 queda acceptada pel que fa a execució completa i memòria
limitada, H5 queda acceptada per l'acceleració amb paritat i H6 queda acceptada
pel cost negligible de la decisió espacial. La Raspberry valida aquests
criteris com a plataforma representativa de baix recurs, però no certifica el
sistema per a vol. El capítol següent resumeix aquestes respostes i delimita el
treball futur.
